\documentclass[11pt]{article}
\usepackage{fontspec}
\usepackage[french]{babel}
\usepackage{amsmath}
\usepackage{amssymb}
\usepackage{amsthm}
\usepackage{mathrsfs}
\usepackage{tikz-cd}
\usepackage{enumitem}
\usetikzlibrary{babel,decorations.markings}
\newcommand{\printedpage}[1]{\marginpar{\footnotesize p.~#1}}
\newcommand{\msp}[1]{\marginpar{\footnotesize [#1]}}
\newcommand{\CN}{\mathcal{C}_{\mathbf{N}^o}}
\tikzset{zigzag/.style={-, postaction={decorate},
  decoration={markings, mark=at position 0.12 with {\arrow{>}},
  mark=at position 0.88 with {\arrow{<}}}}}

\begin{document}

\printedpage{104}
% head: A. Grothendieck : Dérivateurs, Ch. VII, Édition des Universités de Montpellier II et Paris VII
\msp{page 166}

Ainsi, le foncteur $\psi\varphi$ est bien homotope au foncteur identique. On a gagné !

Passons à $1'$, $2'$) pour $\CN$. L'assertion $1'$) se démontre exactement de la même manière que 1) pour \underline{Ch}$(X)$, en regardant $\CN$ comme fibrée sur $\mathbf{N}^o$ (qui a un objet final 0), et en utilisant le diagramme
\[
\begin{tikzcd}
X = (X_{\mathbf{N}^o})_e \arrow[rr, "i'"] & & X_{\mathbf{N}^o} = X \times \mathbf{N}^o \arrow[rr, "j'"] \arrow[dr] & & \mathcal{C}_{\mathbf{N}^o} \arrow[dl] \\
 & & & \mathbf{N}^o \ , &
\end{tikzcd}
\]
où $j'$ est (cartésien et) une Hot-équivalence sur chaque fibre d'après le lemme, donc une Hot-équivalence globalement, et $i'$ est une Hot-équivalence car $\mathbf{N}^o$ est asphérique (ayant un objet final). Donc $j'i'$ est asphérique, q.e.d.

Il reste à prouver que
\[ \CN \longrightarrow X \times X \]
est propre et lisse. Ici encore,
\msp{page 167}
c'est exactement le même argument que tantôt, je présume. Il y a cependant une complication technique. Nous avions pris, pour définir
\[ \underline{t} : \mathbf{N}^o \longrightarrow \underline{\mathrm{Ty}}, \qquad \underline{t} = (t_n)_n, \]
le choix
\[ t_n = (\Delta_+\Delta_-)^{2n}, \]
ou si on préfère
\[ t_n = (\Delta_+\Delta_-)^n, \]
en tous cas les $t_n$ sont à direction initiale $+$, terminale $-$, et ce sont des types réduits. Si on prend $\Delta_+t_n\Delta_-$, on trouve un type non réduit, qui n'est pas parmi les $t_n$. Il ne servirait à rien de changer la définition de $t_n$ en $t_n = (\Delta_-\Delta_+)^{\text{quelque chose}}$, car cette fois $\Delta_+t_n\Delta_-$ sera bien réduit, mais n'aura pas les bons signes initial et terminal. Donc pour définir $\tilde{c}$ en termes de $c$, comme objet de $\CN$, il faudra, si $c$ est de type $\tau = t_n$, prendre $\tilde{c}$ de type $\Delta_+\Delta_-\tau\Delta_+\Delta_- = t_{n+1}$ (en prenant $t_{n+1} = (\Delta_+\Delta_-)^{2n}$), et de même pour définir $\tilde{c}'$ en termes de $c'$. Les définitions pertinentes du passage de $(c',w')$ à $(\tilde{c}',\tilde{w}')$ sont décrites dans la figure : % [sic] t_{n+1} = (\Delta_+\Delta_-)^{2n} -> t_{n+1} = (\Delta_+\Delta_-)^{2n+2}, exponent inconsistent with t_n = (\Delta_+\Delta_-)^{2n}

\msp{page 168}
% foot: Version 22 décembre 2001

\printedpage{105}
% head: A. Grothendieck : Dérivateurs, Ch. VII, Édition des Universités de Montpellier II et Paris VII
\[
\tilde{w}'\Big\downarrow
\left\{
\overbrace{\underbrace{
\begin{tikzcd}[ampersand replacement=\&]
a' \arrow[r, "\mathrm{id}"] \arrow[d, "\mathrm{id}"'] \& a' \arrow[d, "u"'] \& a' \arrow[l, "\mathrm{id}"'] \arrow[d, "u"'] \arrow[rrrr, dotted, no head, "c'", "\longrightarrow"'] \& {} \arrow[d, "w'"] \& {} \arrow[d, "w'"] \& {} \arrow[d, "w'"] \& b' \arrow[r, "\mathrm{id}"] \arrow[d, "v"'] \& b' \arrow[d, "v"] \& b' \arrow[l, "\mathrm{id}"'] \arrow[d, "\mathrm{id}"] \\
a' \arrow[r, "u"] \& a \& a \arrow[l, "\mathrm{id}"'] \arrow[rrrr, no head, "c", "\longrightarrow"'] \& {} \& {} \& {} \& b \arrow[r, "\mathrm{id}"] \& b \& b'. \arrow[l, "v"']
\end{tikzcd}
}_{\tilde{c}}}^{\tilde{c}'}
\right\}
\tilde{w}'
\]

Ainsi
\[
\left\{
\begin{array}{lcl ll}
\tilde{c} & = & \xi_u \circ c \circ \eta_v & \text{où} &
\left\{
\begin{array}{lccccc}
\xi_u : & a' & \xrightarrow{u} & a & \xleftarrow{\mathrm{id}} & a \\[1ex]
\eta_v : & b & \xrightarrow{\mathrm{id}} & b & \xleftarrow{v} & b'
\end{array}
\right\}
\text{ type } \Delta_+\Delta_- \\[4ex]
\tilde{c}' & = & a'_{\Delta_+\Delta_-} \circ c' \circ b'_{\Delta_+\Delta_-} & & \\[4ex]
\tilde{w}' & = & \lambda(u) * w * \mu(v) & \text{avec} & % [sic] w -> w', the homomorphism of the pair (c',w')
\left\{
\begin{array}{lccc}
\lambda(u) : & a'_{\Delta_+\Delta_-} & \longrightarrow & \xi_u \\[1ex]
\lambda(v) : & b'_{\Delta_+\Delta_-} & \longrightarrow & \eta_v % [sic] \lambda(v) -> \mu(v), the formula above has \mu(v)
\end{array}
\right\}
\text{ comme dans \texttt{[la]} figure.}
\end{array}
\right.
\]
On trouve ainsi, mutatis mutandis,
\[ \CN(c;u,v) \xrightarrow{\varphi} \CN(a',b')/\tilde{c}, \qquad \varphi(c',w') = (\tilde{c}',\tilde{w}'), \]
et on définit le foncteur
\[ \CN(c;u,v) \xleftarrow{\psi} \CN(a',b')/\tilde{c} \]
ainsi : à $(c'',w'')$ dans $\CN(a',b')/\tilde{c}$, donc un chemin $c''$ de type $t_n = (\Delta_+\Delta_-)^{2n}$ de $a'$ à $b'$ et un homomorphisme $w'' : c'' \longrightarrow \tilde{c}$ à extrémités fixes, on associe $(c',w')$ dans $\CN(c;u,v)$, avec
\[ c' = c'', \qquad w' = \rho_c w'', \]
où
\[ \rho_c : \tilde{c} \longrightarrow c \]
est défini de façon similaire que tantôt, par
\msp{page 169}
\[
\begin{tikzcd}
a' \arrow[r, "u"] \arrow[d, "u"'] & a \arrow[d, "\mathrm{id}"] & a \arrow[l, "\mathrm{id}"'] \arrow[d, "\mathrm{id}"] \arrow[rr, zigzag] & {} \arrow[d, "\mathrm{id}"'] & b \arrow[r, "\mathrm{id}"] \arrow[d, "\mathrm{id}"] & b \arrow[d, "\mathrm{id}"] & b' \arrow[l, "v"'] \arrow[d, "v"] & \tilde{c},\ \text{ type } t_{n+1}. \\
a \arrow[r, "\mathrm{id}"] & a & a \arrow[l, "\mathrm{id}"'] \arrow[rr, zigzag, "c"] & {} & b \arrow[r, "\mathrm{id}"] & b & b' \arrow[l, "\mathrm{id}"'] & \varphi^*_{n,n+1}(c).
\end{tikzcd}
\]

On doit seulement vérifier que l'endofoncteur $\psi\varphi$ de $\CN(c;u,v)$ est homotope à l'identité. Mais ce foncteur transforme $(c',w')$ en $(\tilde{c}' = a'_{\Delta_+\Delta_-}c'b'_{\Delta_+\Delta_-}, \overline{w'})$, où $\overline{w'}$ est encore le composé
% foot: Version 22 décembre 2001
\printedpage{106}
% head: A. Grothendieck : Dérivateurs, Ch. VII, Édition des Universités de Montpellier II et Paris VII
\[
\begin{tikzcd}
\tilde{c}' \arrow[r, "\text{can.}"] \arrow[rr, bend right=20, "\overline{w'}"'] & c' \arrow[r, "w'"] & c.
\end{tikzcd}
\]

On a donc
\[ \psi\varphi \longrightarrow \mathrm{id}, \]
$\psi\varphi$ est homotope au foncteur identique.

Cela prouve la lissité de $\CN \longrightarrow X \times X$. Pour la propreté, il faut prouver la lissité de le cas de $X^o$, et du système projectif $\underline{t}^{\mathrm{opp}} \simeq (t_n^{\mathrm{opp}})_n$. Mais on a $\underline{t}^{\mathrm{opp}} = \underline{t}$, et il suffit d'appliquer ce qu'on a déjà prouvé au cas de $X^o$, pour voir que $\CN \longrightarrow X \times X$ est propre. % [sic] de le cas -> dans le cas, grammatical

Il ne reste plus que le lemme à prouver. Ça doit être essentiellement
\msp{page 170}
trivial, et je renonce à rédiger une démonstration.

Je vais maintenant montrer `formellement' que la factorisation `canonique' d'une flèche de Cat $f : X' \longrightarrow X$ en
\[ X' \xrightarrow{i_f} \overline{X'} \xrightarrow{\bar{f}} X, \qquad i_f \text{ Hot-équivalence,} \]
avec (on l'espérait) $\bar{f}$ propre et lisse, à partir d'une factorisation similaire pour la flèche diagonale $X \xrightarrow{\ \ \mathrm{diag}_X\ \ } X \times X$ en
\[
\begin{tikzcd}
X \arrow[d, "i_X\ \ \text{\footnotesize Hot-équivalence}"] \\
\mathcal{C}_X\ \ \text{`chemins'} \arrow[d, "p_X"] \\
X \times X
\end{tikzcd}
\]
est chose formelle. (Cf. page 139 pour les notations et \texttt{[la]} problématique.) On regarde le diagramme
\[
\begin{tikzcd}
 & \mathcal{C}_X \arrow[d, "p_X"] \arrow[dl, "\beta_X"'] \arrow[dr, phantom, "\text{cart.}"] & C(f) \arrow[l, "\mathrm{pr}'_2"'] \arrow[d, "p_{X,f}"] \arrow[dll, bend right=50, densely dotted, "\bar{f}"'] & [-1.5em] \simeq X' \times_X (C_X, \sigma = \mathrm{pr}_1 \circ p_X) \\ % ?? uncertain: line style of the curved arrow \bar{f} -- densely dotted / thin dashed / thin solid
X & X \times X \arrow[l, "\mathrm{pr}_2^X"] \arrow[d, "\mathrm{pr}_1^X"'] \arrow[dr, phantom, "\text{cart.}"] & X' \times X \arrow[l, "f \times \mathrm{id}_X = f'"] \arrow[d, "\mathrm{pr}_1"] & \\
 & X & X' \arrow[l, "f"'] \arrow[uu, bend right=60, dotted, "i_f"'] &
\end{tikzcd}
\]
% foot: Version 22 décembre 2001
\printedpage{107}
% head: A. Grothendieck : Dérivateurs, Ch. VII, Édition des Universités de Montpellier II et Paris VII
(\footnote[40]{NB $C(f)$ est défini comme l'image inverse sur $X'$ du fibré $C_X$ sur $X$ (par $\sigma_X = \mathrm{pr}_1^X \circ p_X$) par $f : X' \longrightarrow X$.}). Comme $p_X$ est propre et lisse, i.e. `parfait', son image inverse sur $X' \times X$, savoir $p_{X,f}$, l'est aussi; d'autre part $\mathrm{pr}_2^X \circ f'$ est la
\msp{page 171}
projection $\mathrm{pr}_2$ de $X' \times X$ sur $X$, donc c'est aussi un morphisme parfait (puisque $X' \longrightarrow e$ est parfait, donc il en est de même du composé
\[ (\mathrm{pr}_2^X \underbrace{f') \circ p_{X,f}} = \underbrace{\mathrm{pr}_2^X(p_X}\, \mathrm{pr}'_2) = \beta_X \mathrm{pr}'_2 \overset{\text{déf}}{=} \bar{f} \ \ ). \]

Je rappelle que la section $i_f$ de $C(f)$ sur $X'$ est définie à partir de la section $i_X$ de $(\mathcal{C}_X,\sigma_X)$ sur $X$. C'est une Hot-équivalence, car $\sigma_X = \mathrm{pr}_1^X p_X$ est parfait à fibres asphériques, donc universellement une Hot-équivalence, donc $\mathrm{pr}_1 p_{X,f} : C(f) \longrightarrow X'$ est aussi une Hot-équivalence, et toute section en est donc une aussi.
% foot: Version 22 décembre 2001

\end{document}
